\chapter*{Conclusion générale}
\addcontentsline{toc}{chapter}{Conclusion générale}
\markboth{\textbf{Conclusion générale}}{}

Le service support de \otbs réunit onze consultants au service de huit filiales du groupe. En
quatorze mois, il avait accumulé 1\,800 tickets, dont 476 seulement menés jusqu'au statut
\emph{Résolu}, sans qu'aucun ne soit jamais réutilisé. La question posée en introduction était de
savoir si cet historique pouvait devenir un actif exploitable, sans
qu'aucune donnée ne sorte de l'infrastructure interne et sans dessaisir le consultant de sa
décision.

HelpDesk AI y répond. Une demande déposée en texte libre reçoit une catégorie proposée par un
modèle, accompagnée des tickets déjà traités les plus proches. L'administrateur corrige, traite et
rédige la solution appliquée, laquelle rejoint la base de connaissances interrogée par la demande
suivante. La plateforme compte dix-sept écrans, quatre composants conteneurisés, et deux cent
soixante tests automatisés en protègent le comportement.

Le meilleur système de classification obtenu, un transformer francophone adapté par LoRA et complété
de règles métier, atteint 0,8343 de F1-macro dans la configuration publiée, à seize millièmes de
l'objectif fixé avant l'expérimentation — et 0,8644 lorsqu'il reçoit l'entrée que l'application lui
enverrait réellement, donc au-dessus de cet objectif. La version servie repose sur une chaîne
lexicale plus légère, qui obtient 0,5827 sous ce même protocole, parce que l'inférence du
transformer coûte 128,5\,ms par ticket sur processeur contre 5,0\,ms, à quoi s'ajoute une requête
déjà mesurée au-dessus du budget de latence. Le moteur de recherche hybride rend 2,0 suggestions par
demande, justes à 0,5185 une fois son seuil de pertinence recalibré sur l'échelle réelle de ses
scores — valeur à lire comme majorée, puisque ce seuil a été choisi sur les quarante requêtes qui
l'évaluent, et qu'il perd sept à huit points hors échantillon. Le modèle de langage, enfin,
s'exécute localement sans intervenir sur aucun chemin obligatoire.

Trois exigences chiffrées ne sont pas tenues par le système livré, et nous les énonçons plutôt que
de les laisser découvrir : le temps de réponse dépasse de quelques dizaines de millisecondes le
budget fixé au quatre-vingt-quinzième centile ; le rappel sur la catégorie \emph{Sécurité} et la qualité du classifieur,
visée à 0,85 de F1-macro, ne sont atteints ni par la chaîne servie — 0,5827, et zéro de rappel sur
\emph{Sécurité} au jeu de test gelé — ni par le modèle dans sa configuration publiée. Sur ces deux
dernières, une nuance s'est imposée en cours de route et nous la portons ici : mesuré sur l'entrée
que la production lui fournirait, le transformer franchit les deux seuils, 0,8644 et 0,90. Ce
n'est pas lui qui est servi, l'exigence porte sur ce qui l'est, et elle reste donc non tenue ; mais
l'obstacle n'est plus la qualité du modèle, il est le coût de son exécution et la décision de
l'assumer. La
limite de fond tient au corpus, dont les 1\,800 tickets ne portent que cinquante-huit titres
distincts, si bien que le jeu de test gelé ne couvre que cinq des huit catégories. Trois réserves
subsistent enfin. La confidentialité est tenue sur son volet technique et pas sur l'autre : aucune
donnée ne quitte l'infrastructure, le corpus et la chaîne d'export pseudonymisent les personnes et
un audit de 322 fragments de noms n'en retrouve aucun, mais les durées de conservation attendent un
arbitrage écrit du service et les personnes n'ont pas été informées. L'interface, ensuite, n'est
couverte par aucun test automatique. La valeur d'usage, enfin, n'a pas été mesurée auprès des
consultants.

Servir le transformer redeviendra défendable dès que la machine hôte disposera d'une carte
graphique, où son coût retombe à 58,7\,ms par ticket. La perspective de fond reste l'enrichissement
du corpus par l'exploitation elle-même, chaque ticket traité produisant une étiquette vérifiée :
c'est le seul chemin qui lève la limite de diversité, et le mécanisme est déjà en place.

Ce travail nous a conduits à parcourir les six phases de CRISP-DM, de la compréhension du métier
jusqu'au déploiement. Il nous a surtout appris à distinguer un résultat d'une mesure, et à préférer un
chiffre plus bas qui tient à un chiffre flatteur qui ne tient pas. Un service support ne manque pas
de données, il manque de moyens de s'en servir. La plateforme livrée en fournit un, qui propose sans
décider : elle qualifie, retrouve et suggère, elle ne clôt aucun ticket.
